% Generated draft: requires mathematical and source review.
\documentclass[11pt]{article}
\usepackage[margin=25mm]{geometry}
\usepackage{fontspec}
\setmainfont{Latin Modern Roman}
\usepackage{amsmath,amssymb,mathtools,unicode-math}
\setmathfont{Latin Modern Math}
\usepackage{enumitem,needspace,xurl,hyperref}
\setlength{\parindent}{0pt}
\setlength{\parskip}{5pt}
\setlength{\emergencystretch}{3em}
\newcommand{\N}{\mathbb{N}}
\newcommand{\Z}{\mathbb{Z}}
\newcommand{\Q}{\mathbb{Q}}
\newcommand{\R}{\mathbb{R}}
\newcommand{\C}{\mathbb{C}}
\newcommand{\F}{\mathbb{F}}
\newcommand{\A}{\mathbb{A}}
\newcommand{\PP}{\mathbb{P}}
\newcommand{\E}{\mathbb{E}}
\newcommand{\eps}{\varepsilon}
\newcommand{\dd}{\,\mathrm{d}}
\begin{document}
\section*{UnsolvedMath ID: 30005238}
\textbf{Source problem code:} catalogue entry
\subsection{Definitions and mathematical statement}
The supplied record concerns approximate separability of Helmholtz Green functions in the high-frequency limit. Its displayed equation is
\[
\Delta_x G(x,y)+k^2 n^2(x)G(x,y)=\delta(x-y),\qquad x,y\in\R^d,
\]
where \(\Delta_x\) denotes the Laplacian in the variable \(x\), and \(\delta(x-y)\) is the Dirac distribution at \(y\). Thus the displayed equation is interpreted distributionally in \(x\). The source does not specify the dimension, the admissible wavenumbers \(k\), assumptions on the coefficient \(n\), or a radiation condition selecting a Green function.

A finite separable approximation has the formal form
\[
G(x,y)\approx\sum_{j=1}^{r}a_j(x)b_j(y).
\]
The source discusses the minimum number of terms needed for such an approximation, but its excerpt ends at ``The approximate separability of G(''. It supplies neither the sets on which approximation is required nor an error norm, tolerance, admissible factor spaces, or parameter dependence. Consequently, no precise minimum-rank quantity or quantified asymptotic claim can be recovered.

The supplied excerpt reports lower bounds obtained from an explicit characterization of correlations between Green functions and a tight dimension estimate for the best linear subspace approximating a set of decorrelated Green functions. It also reports upper bounds obtained by constructing separable approximations and sharpness in a few case studies of practical interest. These are descriptions of reported results, without their hypotheses or formulas. The excerpt further describes high separability for Green functions of coercive elliptic differential operators in divergence form, citing an unidentified reference [1], and numerical algorithms exploiting low-rank discretized systems.

\textbf{Unresolved formulation.} The supplied statement-verification field says that the official report presents separability bounds answering a rhetorical title and does not isolate an unresolved mathematical claim; the supplied statement status is \texttt{unrecoverable}. Accordingly, the equation and approximation framework above are the recoverable mathematical content, not a reconstructed open conjecture. The record does not determine a precise additional theorem to prove.
\subsection{Short English statement}
The source studies how many separated terms are needed to approximate a high-frequency Helmholtz Green function and relates this to solution-space complexity and low-rank numerical methods. It reports bounds and sharp examples, but the excerpt is truncated and specifies no precise unresolved claim. Missing coefficient assumptions, radiation conditions, approximation regions, and error criteria prevent a complete mathematical formulation.
\subsection{Sources}
\begin{itemize}
\item Supplied local registry record: permanent UnsolvedMath ID \(30005238\), source problem code \texttt{catalogue entry}, title \emph{High-Frequency Helmholtz Green-Function Separability}, in the Oberwolfach Reports Open Problems collection. No individual proposer is supplied; the proposed year is 2022.
\item Supplied report citation: \emph{At the Interface between Semiclassical Analysis and Numerical Analysis of Wave Scattering Problems} (2023), Oberwolfach Reports, pp.\ 2511--2587, DOI \url{https://doi.org/10.4171/owr/2022/43}. The supplied literature list labels this Oberwolfach Report 43/2022; the excerpt contains the page marker 2529. This is supplied provenance, not an independently inspected source.
\item Supplied literature reference: Engquist--Zhao, CPAM 2018, \url{https://doi.org/10.1002/cpa.21755}. No article title or more detailed bibliographic information is supplied.
\item Literature assessment dated exactly \texttt{2026-08-22}: the supplied curation labels the problem \texttt{partially\_solved} and attributes the difficulty to high-frequency oscillation, pollution, indefiniteness, and Green-function separability. It reports sharp bounds in important settings and says that a complete uniformly sharp theory for general heterogeneous problems is not established. This is supplied dated curation derived from the OWR material, not an independently verified literature assessment; its time-sensitive status and cited claims require independent verification.
\item The excerpt's reference [1], concerning coercive elliptic operators, is not identified in the supplied record. Its identity cannot be inferred from the separate literature list.
\end{itemize}
\end{document}
